% Scattering Angle in the CM System

\subsubsection{Change of variables}
%%%%%%%%%%%%%%%%%%%%%%%%%%%%%%%%%%%
%
The integral in \cref{eq:sn_reduced_theta} is of the form
%
\begin{equation}
    \label{eq:sn_integration_integral_R}
    \int_{R_0}^\infty F(R) \, \mathrm{d}R ~,
\end{equation}
%
where \(R_0\) is a pole of $F(R)$, which is obtained by solving \cref{eq:sn_apsis}, see Section~\ref{s:apsis}.
Having a numerically determined pole of the integrand as an integration boundary is inconvenient, so the first step of calculating the scattering integrals is to change the integration variable.
One possibility is \cite{Eck91}
%
\begin{equation}
    \label{eq:sn_transform1}
    R = \frac{ R_0 }{ (1-u^2) } ~,
\end{equation}
%
which yields
%
\begin{equation}
    \label{eq:sn_integration_integral_u1}
    \int_{R_0}^\infty F(R) \, \mathrm{d}R =
        2 R_0 \int_0^1 F \left( \frac{ R_0 }{ 1-u^2 } \right)
        \frac{ u \, \mathrm{d}u }{ (1-u^2)^2 } ~.
\end{equation}
%
With \cref{eq:sn_reduced_theta}
%
\begin{equation*}
    \pi - \theta = 2 \, P \int_{R_0}^\infty
        \frac{ \mathrm{d}R }{ R^2 \, g(R) }
\end{equation*}
%
we obtain
%
\begin{equation}
    \label{eq:sn_cmangl_legendre_theta_u}
    \pi - \theta = \frac{ 4 P }{ R_0 } \int_0^1
        \frac{ u \, \mathrm{d}u }{ g \left( \frac{ R_0 }{ 1-u^2 } \right) } ~,
\end{equation}
%
where $g(R)$ is given by \cref{eq:sn_reduced_g}
%
\begin{equation*}
    g(R) = \sqrt{ 1 - \frac{ 1 }{ \varepsilon R } \Phi(R)
        - \frac{ P^2 }{ R^2 } } ~.
\end{equation*}
%
In the following we show that the integrand is well-behaved in the interval $[0, 1]$.

\subsubsection{Rewriting the integrand}
%%%%%%%%%%%%%%%%%%%%%%%%%%%%%%%%%%%%%%%
%
It is still tricky to evaluate the integrand of \cref{eq:sn_cmangl_legendre_theta_u} near $u = 0$ (corresponding to $R \approx R_0$), since the evaluation of $g \left ( \frac{ R_0 }{ 1-u^2 } \right)$ heavily depends on the accuracy of the numerically determined value of $R_0$.
To avoid potential problems arising from this we subtract \cref{eq:sn_apsis}
%
\begin{equation*}
    1 - \frac{1}{\varepsilon R_0} \Phi(R_0) - \frac{P^2}{R_0^2} = 0
\end{equation*}

from the expression under the square root of
%
\begin{equation}
    \label{eq:sn_cmangl_legendre_g_u}
    g \left( \frac{ R_0 }{ 1-u^2 } \right) =
    \sqrt{ 1 - \frac{ 1 }{ \varepsilon R_0 }
        \Phi\left( \frac{ R_0 }{ 1-u^2 } \right) (1-u^2)
        - \frac{ P^2 }{ R_0^2 } (1-u^2)^2 } ~,
\end{equation}
%
which yields
%
\begin{equation}
    \label{eq:sn_cmangl_legendre_g_chi}
    g \left( \frac{ R_0 }{ 1-u^2 } \right) =
        \sqrt{ \frac{ R_0 }{ \varepsilon P^2 } \,
               \frac{ \Phi(R_0) - \Phi\left( \frac{ R_0 }{ 1-u^2 } \right) \, (1-u^2) }{ u^2 }
               + ( 2 - u^2 ) } ~.
\end{equation}
%
\Cref{eq:sn_cmangl_legendre_g_chi} can be rewritten
%
\begin{equation}
    \label{eq:sn_cmangl_legendre_g_G}
    g \left( \frac{ R_0 }{ 1-u^2 } \right) =
        u \frac{P}{R_0} \, G(u)
\end{equation}
%
with
%
\begin{equation}
    \label{eq:sn_cmangl_legendre_G_chi}
    G(u) = \sqrt{ \rho \, \chi(u) + (2-u^2) } ~,
\end{equation}
%
~
%
\begin{equation}
    \label{eq:sn_cmangl_legendre_chi}
    \chi(u) = \frac{ \Phi(R_0) - \Phi\left( \frac{ R_0 }{ 1-u^2 } \right) \, (1-u^2) }{ u^2 } ~,
\end{equation}
%
and
%
\begin{equation}
    \label{eq:sn_cmangl_legendre_rho}
    \rho = \frac{ R_0 }{ \varepsilon P^2 } ~.
\end{equation}

\(\chi(u)\) is well-behaved with
%
\begin{equation}
    \label{eq:sn_cmangl_legendre_chi_lim}
    \lim_{u \rightarrow 0} \chi(u) = \Phi(R_0) - R_0 \Phi'(R_0) , \qquad
    \chi(1) = \Phi(R_0) ~.
\end{equation}
%
The former is obtained by l’Hospital’s rule.
Rounding errors can be minimized by replacing \cref{eq:sn_cmangl_legendre_chi} by \cref{eq:sn_cmangl_legendre_chi_lim} for $u$ smaller than the square root of the floating point precision.
It is also easily seen that $\chi(u)$ is always positive for repulsive potentials ($\Phi'(R) < 0$ for $R \ge R_0$).
\(\chi(u)\) is plotted in \cref{fig:sn_cmangl-legendre-chi-plot} for some values of \(R_0\).
%
\begin{figure}[htbp]
    \centering
    \includegraphics[scale=0.5]{chi_ZBL.pdf} ~~
    \includegraphics[scale=0.5]{chi_NLHlin_As_Si.pdf}
    \caption{Function \(\chi(u)\) according to \cref{eq:sn_cmangl_legendre_chi} for the ZBL universal potential (left) and the NLHlin potential for the atom pair As-Si (right).}
    \label{fig:sn_cmangl-legendre-chi-plot}
\end{figure}

\Cref{eq:sn_cmangl_legendre_theta_u} may now be rewritten
%
\begin{equation}
    \label{eq:sn_cmangl_legendre_theta_chi}
    \pi - \theta = 4 \int_0^1  \frac{ 1 }{ G(u) } \, \mathrm{d}u .
\end{equation}
%
The integrand has no poles, and the integral is amenable to standard methods of numerical quadrature.

$G(u)$ evaluates to infinity for $P = 0$. To avoid handling infinite values, we define
%
\begin{equation}
    \label{eq:sn_H_def}
    H(u) = P \, G(u) = \sqrt{\frac{R_0}{\varepsilon} \chi(u) + P^2 (2-u^2)} \, .
\end{equation}
%
Then \cref{eq:sn_cmangl_legendre_theta_chi} can be rewritten
%
\begin{equation}
    \label{eq:sn_cmangl_legendre_theta_H}
    \pi - \theta = 4 P \int_0^1  \frac{ 1 }{ H(u) } \, \mathrm{d}u .
\end{equation}
%
Apart from showing $\theta = \pi$, this explicitly shows that $\pi - \theta \sim P$ for small impact parameters.

\subsubsection{Integration using Gauss-Legendre quadrature}
%%%%%%%%%%%%%%%%%%%%%%%%%%%%%%%%%%%%%%%%%%%%%%%%%%%%%%%%%%%
%
Using Gauss-Legendre quadrature, we can approximate \cref{eq:sn_cmangl_legendre_theta_chi} by
%
\begin{equation}
    \label{eq:sn_cmangl_legendre_theta_chi_approx_plain}
    \pi - \theta \approx
        4 \sum_{i=1}^n \frac{ \tilde{w}_i }{ G(\tilde{u}_i) }
\end{equation}
%
with \(\tilde{u}_i\) and \(\tilde{w}_i\) from \cref{eq:sn_legendre_transform}. With $a=0$ and $b=1$ they read
%
\begin{equation}
    \label{eq:sn_gauss_legendre_transform_infinite}
    \tilde{u}_i = \frac{ u_i + 1 }{ 2 }, \qquad
    \tilde{w} = \frac{ w_i }{ 2 }.
\end{equation}

\Cref{eq:sn_cmangl_legendre_theta_chi_approx_plain}, however, has the problem that $\theta$ does not tend to zero exactly as $P \rightarrow \infty$ or $\varepsilon \rightarrow \infty$, since the right-hand side only approximates $\pi$.
Note that for \(\rho \rightarrow 0\) the integral in \cref{eq:sn_cmangl_legendre_theta_chi} becomes
%
\begin{equation}
    \label{eq:sn_cmangl_legendre_pi}
    \int_0^1 \frac{ 1 }{ \sqrt{ 2-u^2 } } \, \mathrm{d}u = \frac{ \pi }{ 4 } ~,
\end{equation}
%
and \(\theta \rightarrow 0\) as it should.
Since the integrand is not a polynomial, the Legendre approximation of the integral will not exactly evaluate to \(\pi/4\), and \cref{eq:sn_cmangl_legendre_theta_chi_approx_plain} will not tend to $\pi$ as \(\rho \rightarrow 0\).
This may be fixed by scaling the integral to the correct value of \(\pi/4\):
%
\begin{equation}
    \label{eq:sn_cmangl_legendre_theta_chi_approx}
    \pi - \theta \approx
        \frac{ \pi }{ \pi_\mathrm{num}/4 }
        \sum_{i=1}^n \frac{ \tilde{w}_i }{ G(\tilde{u}_i) }
\end{equation}
%
with
%
\begin{equation}
    \label{eq:sn_cmangl_legendre_pinum}
    \frac{\pi_\mathrm{num}}{4} =
        \sum_{i=1}^n \frac{ \tilde{w}_i  }{ \sqrt{ 2-\tilde{u}_i^2 } }
\end{equation}
%
Since the sum in \cref{eq:sn_cmangl_legendre_theta_chi_approx} evaluates to \(\pi_\mathrm{num}/4\) for \(\rho \rightarrow 0\), $\pi - \theta$ tends to $\pi$ and \(\theta \rightarrow 0\).
At the same time, for \(\rho \rightarrow \infty\) (i.e., for $\varepsilon \rightarrow 0$ or $P \rightarrow 0$), $G(u) \rightarrow \infty$, and the sum in \cref{eq:sn_cmangl_legendre_theta_chi_approx} goes to zero and \(\theta \rightarrow \pi\) as it should.

In terms of the function $H(u)$, \cref{eq:sn_H_def}, \cref{eq:sn_cmangl_legendre_theta_chi_approx} reads
%
\begin{equation}
    \label{eq:sn_cmangl_legendre_theta_H_approx}
    \pi - \theta \approx
        \frac{ \pi }{ \pi_\mathrm{num}/4 }  \,
        P \sum_{i=1}^n \frac{ \tilde{w}_i }{ H(\tilde{u}_i) }
\end{equation}

One can also apply this correction to the Gauss-Legendre approximation of \cref{eq:sn_cmangl_legendre_theta_u}:
%
\begin{equation*}
    %\label{eq:sn_cmangl_legendre_theta_g_approx}
    \pi - \theta \approx
        \frac{ \pi }{ \pi_\mathrm{num}/4 } \, \frac{ P }{ R_0 }
        \sum_{i=1}^n
        \frac{ \tilde{w}_i \tilde{u}_i
        }{ g \left( \frac{ R_0 }{ 1-\tilde{u}_i^2 } \right) } \, .
\end{equation*}
%
This expression is not numerically equivalent to \cref{eq:sn_cmangl_legendre_theta_chi_approx}.
It does not require the evaluation of \(\Phi(R_0)\), which might represent a noticable saving of computation time if few abscissas are used.
For few abscissas, the results are very similar.
With increasing number of abscissas, the results are increasingly dependent on the accuracy of the apsis \(R_0\).
The term under the square root of $g(\frac{R_0}{1-u^2})$, \cref{eq:sn_cmangl_legendre_g_chi}, or $G(u)$, \cref{eq:sn_cmangl_legendre_G_chi}, may even become negative, if the accuracy of \(R_0\) is low and the number of abscissas is large.

\subsubsection{Modifications for finite-range screened Coulomb potentials}
%%%%%%%%%%%%%%%%%%%%%%%%%%%%%%%%%%%%%%%%%%%%%%%%%%%%%%%%%%%%%%%%%%%%%%%%%
%
If $\Phi(R) = 0$ for $R$ larger than some $R_\mathrm{max}$,
%
\begin{equation}
    \Phi(R) = 0 \quad \text{for} \quad R \ge R_\mathrm{max}
\end{equation}
%
then for $P \ge R_\mathrm{max}$ we trivially have $\theta = 0$.
Otherwise the integral in \cref{eq:sn_reduced_theta} must be divided in two parts:
%
\begin{equation}
    \label{eq:sn_reduced_theta_finite}
    \pi - \theta =
        2 \, P \int_{R_0}^\infty \frac{ \mathrm{d}R }{ R^2 \, g(R) } =
        2 \, P \int_{R_0}^{R_\mathrm{max}} \frac{ \mathrm{d}R }{ R^2 \, g(R) } +
        2 \, P \int_{R_\mathrm{max}}^{\infty} \frac{ \mathrm{d}R }{ R^2 \, g(R) } ~.
\end{equation}
%
The second term on the right-hand side may be evaluated analytically using $g(R)$ with $\Phi(R) = 0$
%
\begin{equation}
    \label{eq:sn_g_phi=0}
    g(R) = \sqrt{ 1 - \frac{ P^2 }{ R_0^2 }}
\end{equation}
%
and \cite[2.275.4]{Gra80}
%
\begin{equation}
    \label{eq:sn_reduced_theta_finite_2}
    2 \, P \int_{R_\mathrm{max}}^{\infty} \frac{ \mathrm{d}R }{ R^2 \, g(R) } =
        2 \, P \int_{R_\mathrm{max}}^{\infty}
            \frac{ \mathrm{d}R }{ R^2 \sqrt{ 1 - \frac{P^2}{R^2}} } =
        2 \arcsin \frac{P}{R_\mathrm{max}} ~.
\end{equation}
%
Note that for $R_\mathrm{max} \rightarrow \infty$, this term vanishes.

Applying the substitution \cref{eq:sn_transform1} to the remaining integral yields
%
\begin{equation}
    \pi - \theta =
        4 \int_0^{u_\mathrm{max} } \frac{ \mathrm{d}u }{G(u)} +
        2 \arcsin \frac{P}{R_\mathrm{max}}
\end{equation}
%
with
%
\begin{equation}
    \label{eq:sn_umax}
    u_\mathrm{max} = \sqrt{ 1 - \frac{ R_0 }{ R_\mathrm{max} } } ~,
\end{equation}

The integral can be evaluated by Gauss-Legendre quadrature
%
\begin{equation}
    \label{eq:sn_cmangl_legendre_theta_chi_finite_approx_plain}
    \pi - \theta =
        4 \sum_{i=1}^n \frac{ \tilde{w}_i }{G(\tilde{u}_i)} +
        2 \arcsin \frac{P}{R_\mathrm{max}}
\end{equation}
%
with
%
\begin{equation}
    \label{eq:sn_gauss_legendre_transform_finite}
    \tilde{u}_i = u_\mathrm{max} \frac{ u_i + 1 }{ 2 }, \qquad
    \tilde{w}_i = u_\mathrm{max} \frac{ w_i }{ 2 }
\end{equation}
%
which follows from \cref{eq:sn_legendre_transform} with $a=0$ and $b=u_\mathrm{max}$.

As with \cref{eq:sn_cmangl_legendre_theta_chi_approx_plain}, for numerical reasons $\rho \rightarrow 0$ might not exactly result in $\theta \rightarrow 0$ as it should.
Unlike for infinite-range potentials, this does not happen for $P \rightarrow \infty$, because for $P > R_\mathrm{max}$ we have the trivial solution $\theta = 0$.
Also for $P = R_\mathrm{max}$ there is no problem because the first term evaluates to zero and the second term evaluates to $\pi$.
But for very high energies $\varepsilon \rightarrow \infty$ and not too small $P$, $\rho$ can become very small.
In that case,
%
\begin{equation}
    \lim_{\varepsilon \rightarrow \infty} \left[
        4 \int_{0}^{u_\mathrm{max}} \frac{ \mathrm{d} u}{G(u)} \right] =
        4 \int_{0}^{u_\mathrm{max}} \frac{ \mathrm{d} u}{\sqrt{2-u^2}} =
        4 \arcsin \frac{u_\mathrm{max}}{\sqrt{2}} = 2 \arccos \frac{R_0}{R_\mathrm{max}}
\end{equation}
%
One can therefore numerically correct \cref{eq:sn_cmangl_legendre_theta_chi_finite_approx_plain} by
%
\begin{equation}
    \label{eq:sn_cmangl_legendre_theta_chi_finite}
    \pi - \theta =
        2 \, \frac{\arccos (R_0 / R_\mathrm{max})}{\frac{1}{2} \arccos_\mathrm{num}} \,
            \sum_{i=1}^n \frac{ \tilde{w}_i }{G( \tilde{u}_i)}
        + 2 \arcsin \frac{P}{R_\mathrm{max}}
\end{equation}
%
with
%
\begin{equation}
    \frac{1}{2} \arccos_\mathrm{num} =
        \sum_{i=1}^n \frac{ \tilde{w}_i }{\sqrt{2 - \tilde{u}_i^2}} \, .
\end{equation}
$\theta$ tends to 0 as $\varepsilon \rightarrow \infty$, since $\rho \rightarrow 0$ and $R_0 \rightarrow P$, and $\arccos x + \arcsin x = \pi/2$. Contrary to $\pi_\mathrm{num}$, \cref{eq:sn_cmangl_legendre_pinum}, the second sum in \cref{eq:sn_cmangl_legendre_theta_chi_finite} depends on $\varepsilon$ and $P$ via $u_\mathrm{max}$ and $R_0$.

For $P = 0$, $\rho \rightarrow \infty$ and therefore $G(u) \rightarrow \infty$, and both terms vanish. So there is no problem in evaluating \cref{eq:sn_cmangl_legendre_theta_chi_finite} for head-on collisions.
Contrary to $\pi_\mathrm{num}/4$, \cref{eq:sn_cmangl_legendre_pinum}, $\frac{1}{2} \arccos_\mathrm{num}$ depends on $\varepsilon$ and $P$ via $u_\mathrm{max}$ and $R_0$.

In terms of the function $H(u)$, \cref{eq:sn_H_def}, \cref{eq:sn_cmangl_legendre_theta_chi_finite} reads
%
\begin{equation}
    \label{eq:sn_cmangl_legendre_theta_H_finite}
    \pi - \theta =
        2 \, \frac{\arccos (R_0 / R_\mathrm{max})}{\frac{1}{2}\arccos_\mathrm{num}} \,
        P \sum_{i=1}^n \frac{ \tilde{w}_i }{H( \tilde{u}_i)}
        + 2 \arcsin \frac{P}{R_\mathrm{max}} \, .
\end{equation}
%
Also for finite-range interatomic potentials, $\pi - \theta \sim P$ for small P, since $\arcsin x \approx x$ for small $x$.

\subsubsection{Convergence}
%%%%%%%%%%%%%%%%%%%%%%%%%%%
%
\begin{figure}[htbp]
    \centering
    \includegraphics[scale=0.5]{error_legendre_theta_eps1e-4_zbl.pdf} ~~
    \includegraphics[scale=0.5]{error_legendre_theta_eps1e-4_nlhlin.pdf} \\ \smallskip
    \includegraphics[scale=0.5]{error_legendre_theta_eps0.1_zbl.pdf} ~~
    \includegraphics[scale=0.5]{error_legendre_theta_eps0.1_nlhlin.pdf} \\ \smallskip
    \includegraphics[scale=0.5]{error_legendre_theta_eps100_zbl.pdf} ~~
    \includegraphics[scale=0.5]{error_legendre_theta_eps100_nlhlin.pdf}
    \caption{Relative errors in \(\theta\) of n-point Gauss-Legendre quadrature as a function of reduced impact parameter for three energies using \cref{eq:sn_cmangl_legendre_theta_chi_finite} with the number of abscissae $n=1$,2, 3, 4, 5, and 10.
        As interatomic potentials the universal ZBL potential (left) and the NLHlin potential for $Z_1=33$ and $Z_2=14$ were used together with a screening length given by \cref{eq:sn_aNLH}.
        For the ZBL potential the results of Biersack's magic formula \cite{Bie80} (check!) are also included with dashed lines.
        The reference simulations were evaluated using \cref{eq:sn_cmangl_legendre_theta_chi_finite} with 32-point Gauss-Legendre quadrature.
        The relative accuracy of the apsis is better than \(10^{-3}\).}
    \label{fig:error-legendre-theta}
\end{figure}
%
In \cref{fig:error-legendre-theta}, the relative errors in \(\theta\) as calculated by \cref{eq:sn_cmangl_legendre_theta_chi_finite} are shown as a function of reduced impact parameter $P$ for reduced energies \(\varepsilon = 10^{-4}\), \(\varepsilon = 10^{-1}\), and \(\varepsilon = 100\), and different numbers of abscissae $n$ in the Gauss-Legendre quadrature.
The results for the ZBL potential are shown in the left column, while those for the NLHlin potential are found in the right column.
For the ZBL potential, the result of Biersacks magic formula \cite{Bie80} (check!) are also shown.
Notably, the accuracy of the magic formula, about 1~\%, is already reached with only $n=2$ abscissae in the Gauss-Legendre quadrature.
An accuracy 10$^{-4}$--10$^{-3}$ is reached with $n = 4$.
For the NLHlin potential, convergence is even faster in general.

\textbf{TODO: Magic formula}